\documentclass[12pt,a4paper]{article}
\usepackage[utf8]{inputenc}
\usepackage[russian]{babel}
\usepackage{geometry}
\usepackage{graphicx}
\usepackage{listings}
\usepackage{xcolor}
\usepackage{hyperref}
\usepackage{float}

\geometry{left=2cm,right=2cm,top=2cm,bottom=2cm}

\lstset{
    language=Python,
    basicstyle=\footnotesize\ttfamily,
    keywordstyle=\color{blue},
    commentstyle=\color{green!40!black},
    stringstyle=\color{red},
    showspaces=false,
    showstringspaces=false,
    numbers=left,
    numberstyle=\tiny,
    breaklines=true
}

\title{
    \textbf{Лабораторная работа №6} \\
    \large{Программа "Солнечная система"} \\
    \textbf{Групповой отчёт}
}

\author{
    \textbf{Лидер:} Руш Егор \\
    \textbf{Кодировщик:} Руш Егор \\
    \textbf{Тестировщик:} Руш Егор
}

\date{\today}

\begin{document}

\maketitle

\tableofcontents
\newpage

\section{Введение}

Данная программа представляет собой информационную систему для визуализации Солнечной системы. Программа разработана на языке Python с использованием библиотеки Tkinter.

\subsection{Цель работы}
Разработать кросплатформенное приложение для визуализации Солнечной системы с возможностью получения информации о планетах.

\subsection{Задачи}
\begin{enumerate}
    \item Создать графический интерфейс с использованием Tkinter
    \item Реализовать моделирование движения планет по орбитам
    \item Реализовать отображение информации о планетах
    \item Добавить регулировку скорости движения
    \item Обеспечить соответствие PEP8
\end{enumerate}

\section{Архитектура программы}

\subsection{Диаграмма классов}

Программа состоит из следующих основных классов:

\begin{itemize}
    \item \textbf{Планета} - класс, представляющий планету
    \item \textbf{СолнечнаяСистема} - класс, управляющий планетами
    \item \textbf{СолнечнаяСистемаGUI} - класс графического интерфейса
\end{itemize}

\subsection{Структура проекта}

\begin{verbatim}
solar-system/
├── main.py          # Главный файл
├── planets.py       # Классы планет
├── gui.py           # Графический интерфейс
├── data.py          # Данные о планетах
├── utils.py         # Вспомогательные функции
└── README.md        # Документация
\end{verbatim}

\section{Описание классов}

\subsection{Класс Планета}

Класс представляет планету в Солнечной системе.

\textbf{Атрибуты:}
\begin{itemize}
    \item имя - название планеты
    \item радиус\_орбиты - расстояние от Солнца
    \item физический\_радиус - размер планеты
    \item цвет - цвет планеты
    \item угол - угол на орбите
    \item период - период обращения
    \item скорость - угловая скорость
\end{itemize}

\textbf{Методы:}
\begin{itemize}
    \item обновить\_позицию() - обновляет положение
    \item получить\_координаты() - возвращает координаты
    \item установить\_скорость() - устанавливает скорость
\end{itemize}

\subsection{Класс СолнечнаяСистема}

Класс управляет всеми планетами.

\textbf{Атрибуты:}
\begin{itemize}
    \item планеты - список планет
    \item скорость - множитель скорости
\end{itemize}

\textbf{Методы:}
\begin{itemize}
    \item обновить() - обновляет все планеты
    \item установить\_скорость() - устанавливает скорость
    \item найти\_планету\_по\_координатам() - поиск планеты
\end{itemize}

\subsection{Класс СолнечнаяСистемаGUI}

Класс графического интерфейса.

\textbf{Атрибуты:}
\begin{itemize}
    \item root - корневое окно
    \item canvas - холст для рисования
    \item информация\_текст - текстовое поле с информацией
    \item ползунок\_скорости - ползунок регулировки
\end{itemize}

\section{Интерфейс пользователя}

\subsection{Главное окно}

Главное окно программы разделено на три области:
\begin{enumerate}
    \item Центральная область - холст с Солнечной системой
    \item Правая верхняя область - информация о планете
    \item Правая нижняя область - управление скоростью
\end{enumerate}

\subsection{Управление}

\begin{itemize}
    \item Клик по планете - отображение информации
    \item Клик по Солнцу - информация о Солнце
    \item Ползунок скорости - изменение скорости
    \item Кнопка "Сбросить скорость" - возврат к 1.0
\end{itemize}

\section{Тестирование}

\subsection{План тестирования}

\begin{enumerate}
    \item Проверка запуска программы
    \item Проверка отображения всех планет
    \item Проверка движения планет
    \item Проверка отображения информации
    \item Проверка регулировки скорости
\end{enumerate}

\subsection{Результаты тестирования}

Все тесты успешно пройдены. Программа работает стабильно.

\subsection{Ошибки и их решение}

\begin{tabular}{|p{5cm}|p{5cm}|p{5cm}|}
\hline
\textbf{Ошибка} & \textbf{Причина} & \textbf{Решение} \\
\hline
ModuleNotFoundError & Отсутствует модуль tkinter & Установить Python с tkinter \\
\hline
Не отображаются планеты & Неправильные координаты & Проверить центр системы \\
\hline
Не работает ползунок & Неверная привязка события & Проверить command \\
\hline
\end{tabular}

\section{Заключение}

В ходе работы была разработана программа для визуализации Солнечной системы с возможностью получения информации о планетах. Программа соответствует всем требованиям.

\subsection{Выводы}
\begin{itemize}
    \item Программа успешно визуализирует Солнечную систему
    \item Реализовано движение планет с регулируемой скоростью
    \item При нажатии на планету отображается полная информация
    \item Код соответствует стандарту PEP8
    \item Интерфейс интуитивно понятен пользователю
\end{itemize}

\subsection{Возможные улучшения}
\begin{itemize}
    \item Добавить возможность масштабирования
    \item Добавить спутники планет
    \item Добавить анимацию вращения планет
    \item Добавить возможность добавления новых планет
\end{itemize}

\section{Список литературы}

\begin{enumerate}
    \item Документация Python: \url{https://docs.python.org/}
    \item Документация Tkinter: \url{https://docs.python.org/3/library/tkinter.html}
    \item Данные о планетах: \url{https://nssdc.gsfc.nasa.gov/planetary/}
    \item PEP 8 -- Style Guide: \url{https://www.python.org/dev/peps/pep-0008/}
\end{enumerate}

\section{Приложение}

\subsection{Код программы}

Весь код программы представлен в файлах проекта.

\end{document}